\chapter{Sistemas de Posicionamiento de Electrodos \acrshort{eeg}}
\label{chap:anexo_sistemas_electrodos}

\section{Introducción a los Sistemas de Posicionamiento}
\label{sec:intro_sistemas_posicionamiento}

Los sistemas de posicionamiento de electrodos \acrshort{eeg} proporcionan métodos estandarizados para la colocación reproducible de sensores sobre el cuero cabelludo. Estos sistemas garantizan que las mediciones sean comparables entre diferentes laboratorios, investigadores y sesiones experimentales, siendo fundamentales para la validez científica y la reproducibilidad en neurociencia.

La estandarización del posicionamiento de electrodos es crucial porque la actividad eléctrica cerebral presenta patrones espaciales específicos que deben capturarse de forma consistente. Los sistemas más utilizados son el sistema internacional 10-20 y sus extensiones modernas, que se basan en medidas anatómicas precisas del cráneo para determinar las ubicaciones óptimas de los electrodos.

\section{Sistema Internacional 10-20}
\label{sec:sistema_10_20}

El \gls{sistema_10_20}, establecido por la Federación Internacional de Neurofisiología Clínica~\cite{jasper1958report,klem1999ten}, constituye el estándar fundamental para el posicionamiento de electrodos \acrshort{eeg}. Como se muestra en la Figura~\ref{fig:10_20_system}, este sistema utiliza puntos anatómicos de referencia específicos del cráneo para establecer ubicaciones precisas y reproducibles de los electrodos.

\begin{figure}[hp!]
\centering
\includegraphics[width=0.9\textwidth,height=15cm,keepaspectratio]{imaxes/10_20_system.png}
\caption{Ubicaciones de electrodos del sistema internacional 10-20 para registro electroencefalográfico~\cite{tomaton2010electrode}}
\label{fig:10_20_system}
\end{figure}

El sistema se denomina "10-20" porque utiliza distancias del 10\% y 20\% entre puntos anatómicos específicos del cráneo, incluyendo el \textit{nasion} (punto entre la frente y la nariz), el \textit{inion} (protuberancia occipital externa), y los puntos preauriculares (delante de cada oreja). Esta metodología garantiza que las posiciones de los electrodos se adapten proporcionalmente al tamaño del cráneo de cada individuo, manteniendo la consistencia relativa entre diferentes sujetos.

\section{Sistema 10-10 con Nomenclatura Combinatorial Modificada}
\label{sec:sistema_10_10_mcn}

La extensión moderna del \gls{sistema_10_20} incluye el \gls{sistema_10_10}, que proporciona mayor resolución espacial mediante la subdivisión de las distancias anatómicas originales. La Figura~\ref{fig:mcn_system} ilustra este sistema extendido utilizando la nomenclatura combinatorial modificada (\acrshort{mcn}), que permite identificar de forma unívoca cada posición de electrodo mediante una codificación determinada.

\begin{figure}[hp!]
\centering
\includegraphics[width=0.9\textwidth,height=15cm,keepaspectratio]{imaxes/mcn_system.png}
\caption{Posiciones de electrodos \acrshort{eeg} en el \gls{sistema_10_10} utilizando nomenclatura combinatorial modificada, junto con los puntos fiduciales y lóbulos cerebrales asociados~\cite{krol2020eeg}}
\label{fig:mcn_system}
\end{figure}

El \gls{sistema_10_10} con \acrshort{mcn} presenta ventajas significativas para aplicaciones modernas de \acrshort{eeg}, incluyendo mayor densidad de electrodos para mejor resolución espacial, nomenclatura que facilita la identificación automática de posiciones, y compatibilidad con sistemas de localización de fuentes que requieren información espacial precisa. Esta extensión es especialmente relevante para aplicaciones \acrshort{bci} avanzadas que requieren captura de alta resolución espacial.

\section{Aplicación en Sistemas \acrshort{bci}}
\label{sec:aplicacion_sistemas_bci}

Los sistemas de posicionamiento estandarizados son fundamentales para el desarrollo de aplicaciones \acrshort{bci} reproducibles y comparables. La selección de posiciones específicas de electrodos debe basarse en el conocimiento neurofisiológico de las regiones cerebrales relevantes para cada aplicación. Por ejemplo, las aplicaciones basadas en \gls{imaginacion_motora} típicamente utilizan electrodos sobre las áreas sensoriomotoras (C3, C4, Cz), mientras que los \gls{paradigmas_potenciales_evocados} pueden requerir configuraciones más extensas.

La implementación correcta de estos sistemas de posicionamiento en herramientas de desarrollo \acrshort{bci} facilita la integración con bibliotecas de análisis estándar como \textbf{MNE-Python}~\cite{gramfort2013meg} y \textbf{EEGLAB}~\cite{delorme2004eeglab}, que incorporan estas convenciones de nomenclatura y posicionamiento en sus funcionalidades de procesamiento y visualización.